\subsection{General abelian index}

\begin{theorem}[General carrier formula for the abelian coefficient]
For a branch with $N_{\mathrm{fam}}$ copies of the carrier family and $N_H$ light weak doublets from $E_2$, one has
\[
b_Y(N_{\mathrm{fam}},N_H)
:= \frac23 N_{\mathrm{fam}}\tr_{S^+}Y^2 + \frac13 N_H\tr_{E_2}Y^2.
\]
For the minimal carrier this becomes
\[
b_Y(N_{\mathrm{fam}},N_H)
= \frac23 N_{\mathrm{fam}}\cdot\frac{10}{3} + \frac13 N_H\cdot\frac12
= \frac{20}{9}N_{\mathrm{fam}}+\frac16 N_H.
\]
Therefore
\[
\bone(N_{\mathrm{fam}},N_H)
= \frac35\,b_Y(N_{\mathrm{fam}},N_H)
= \frac43 N_{\mathrm{fam}}+\frac1{10}N_H.
\]
\end{theorem}

\begin{corollary}[Canonical carrier branch]
For the canonical branch $(N_{\mathrm{fam}},N_H)=(3,1)$,
\[
b_Y=\frac{41}{6},
\qquad
\bone=\frac{41}{10}.
\]
\end{corollary}

\begin{remark}
The symbol $\nu^c$ is already part of the hard packet $S^+$. Removing it does not define a different carrier; it is an external truncation of the family packet and is treated as such only in Section~7.
\end{remark}

\section{Boundary Multiplicity and Admissible Occupancy}

\subsection{Family-index bridge}

\begin{conjecture}[Family-index bridge]
Let each distinguished cycle $C_i$ carry an integer family index $I_{\mathrm{fam}}(C_i)$. Then the v2.8 block coefficient may be organized as
\[
k_B = \frac12\sum_{i=1}^{3} I_{\mathrm{fam}}(C_i)\,I_1(B).
\]
\end{conjecture}

Combined with the v2.8 block formula
\[
k_B=\frac32\,I_1(B),
\]
this gives
\[
\sum_{i=1}^{3} I_{\mathrm{fam}}(C_i)=3.
\]
The equal-one assignment
\[
I_{\mathrm{fam}}(C_1)=I_{\mathrm{fam}}(C_2)=I_{\mathrm{fam}}(C_T)=1
\]
is the minimal symmetric branch and yields
\[
N_{\mathrm{fam}}=3.
\]

\subsection{Admissible occupancy and topological surcharge}

\begin{definition}[Admissible occupancy]
The admissible chiral occupancy is
\[
\Oadm := \tr_{\Hch}\mathbf 1 = N_{\mathrm{fam}}\dim S^+.
\]
\end{definition}

On the canonical branch one therefore has
\[
\Oadm = 3\times 16 = 48.
\]

\begin{conjecture}[Occupancy transgression]
The topological surcharge should be written as
\[
\deltatop = \Oadm\,\cthree^4.
\]
\end{conjecture}

For the canonical branch this becomes
\[
\deltatop = 48\cthree^4 = \frac{3}{256\pi^4} = 3\cdot 16\,\cthree^4.
\]
In this reading the factor $3$ comes from the weighted boundary geometry while the factor $16$ comes from the minimal carrier packet itself.

\subsection{Seam descent normalization of \texorpdfstring{$\cthree$}{c3}}

The normalization of $\cthree$ is best presented in two stages.

\textbf{Stage 1: hard seam geometry.}
On the spin lift of the minimal seam loop one has
\[
\widetilde U_\Gamma(2\pi)=-1,
\qquad
\widetilde U_\Gamma(4\pi)=1.
\]
So the spin-lifted loop closes only after $4\pi$.

\textbf{Stage 2: conditional axionic identification.}
If the admissible carrier modulus is identified with the axionic angle variable, write
\[
a=\frac{\theta}{4\pi}.
\]
Then
\[
\mathcal L_\theta = \frac{\theta}{32\pi^2}F\widetilde F
= \frac{1}{8\pi}aF\widetilde F,
\]
hence
\[
\cthree = \frac{1}{8\pi}.
\]
The geometry and the final normalization are thus explicitly separated: the $4\pi$ closure is geometric, while the identification of the carrier modulus with the axionic angle remains the extra conditional step.

\section{Electromagnetic Stationarity and the Carrier CFE}

\subsection{Carrier form of the stationary equation}

The electromagnetic closure channel should first be written on a general admissible branch. Define
\[
\phiz(\alpha)
= \frac{1}{6\pi}
+ \Oadm\,\cthree^4 e^{-2\alpha}
+ B\gamma\,\Oadm^2\cthree^8 e^{-4\alpha}
+ \cdots.
\]
Then the carrier form of the CFE is
\[
\alpha^3
-2\cthree^3\alpha^2
-8\cthree^6\,\bone(N_{\mathrm{fam}},N_H)
\ln\!\bigl(\phiz(\alpha)^{-1}\bigr)
=0.
\]
This is the preferred v2.9 presentation because the coefficients are now grouped by carrier data rather than as parallel numerical motifs.

\begin{proposition}[Canonical carrier branch]
For
\[
(N_{\mathrm{fam}},N_H,\nu^c)=(3,1,\text{present}),
\]
one has
\[
\Oadm=48,
\qquad
\bone=\frac{41}{10},
\qquad
B\gamma=\frac32\cdot\frac56=\frac54.
\]
Hence
\[
\phiz(\alpha)
= \frac{1}{6\pi}
+ 48\cthree^4 e^{-2\alpha}
+ \frac54\,48^2\cthree^8 e^{-4\alpha}
+ \cdots.
\]
\end{proposition}

\begin{remark}
The identity
\[
\frac54 = B\gamma
\]
is the compressed carrier reading of the v2.8 two-defect weight.
\end{remark}

Using
\[
\cthree=\frac{1}{8\pi},
\qquad
\phitree=\frac{1}{6\pi},
\qquad
\deltatop=48\cthree^4,
\]
the positive stationary root remains
\[
\alpha^{-1}(0)=\num{137.03599921615843}.
\]
The point of v2.9 is therefore not to move the number, but to reduce the number of primitive inputs behind it.

\begin{table}[H]
\centering
\small
\begin{tabularx}{\linewidth}{@{}lXl@{}}
\toprule
Quantity & Carrier meaning & Value \\
\midrule
$\cthree$ & canonical seam-descent normalization & $\frac{1}{8\pi}$ \\
$\Oadm$ & admissible chiral occupancy & $48$ \\
$\gamma$ & carrier norm $\tr_EY^2$ & $\frac56$ \\
$B$ & carrier rank ratio $\rank E_3/\rank E_2$ & $\frac32$ \\
$B\gamma$ & compressed two-defect weight & $\frac54$ \\
$\bone$ & canonical abelian coefficient & $\frac{41}{10}$ \\
$\alpha^{-1}(0)$ & positive root of the carrier CFE & \num{137.03599921615843} \\
\bottomrule
\end{tabularx}
\caption{Numerical summary of the canonical carrier branch.}
\end{table}

\subsection{Algebraic corollaries of the carrier kernel}

Evaluated on the stationary root of the canonical branch, several previously separate formulas become algebraic corollaries of the same carrier package:
\[
g_{a\gamma\gamma}=-4\cthree=-\frac{1}{2\pi},
\qquad
\beta_{\mathrm{rad}}=\frac{\phiz}{4\pi},
\]
\[
\lambda_C=\sqrt{\phiz(1-\phiz)},
\qquad
\sin^2\theta_{13}=\phiz e^{-\gamma}.
\]
The structural gain is not that these quantities are numerically new, but that they are now visibly downstream of $\Cmin$ and the occupancy branch.

\subsection{Conditional \texorpdfstring{$\chi$}{chi}-channel corollaries}

If the $\chi$ channel and the induced Einstein channel are successfully completed, then the inflationary block can be written as
\[
\frac{M}{\chi_0}=\frac{2}{\sqrt{\pi}}\,\cthree^4,
\qquad
\frac{M}{\Mpl}=\sqrt{8\pi}\,\cthree^4,
\]
and therefore
\[
A_s=\frac{N^2}{24\pi^2}\left(\frac{M}{\Mpl}\right)^2
= \frac{N^2}{3\pi}\cthree^8.
\]
For $N=56$ this yields
\[
A_s=\num{2.09019136432946e-9}.
\]
These results are kept separate from the hard carrier corollaries because they depend on the success of the relative-$\chi$ and gravity channels.

\section{Discrete Exclusion Structure}

The carrier-centered rewrite is strongest when it excludes wrong branches, not when it celebrates an isolated number. The relevant deformations fall into three logical classes.

\begin{table}[H]
\centering
\small
\begin{tabularx}{\linewidth}{@{}l l X l@{}}
\toprule
Logical type & Deformation & Carrier/CFE consequence & Result \\
\midrule
Boundary multiplicity & $N_{\mathrm{fam}}=2$ & $\Oadm=32$ and $\bone=\frac{83}{30}$ in the carrier CFE & $\alpha^{-1}(0)=\num{156.094045457}$ \\
Boundary multiplicity & $N_{\mathrm{fam}}=4$ & $\Oadm=64$ and $\bone=\frac{163}{30}$ in the carrier CFE & $\alpha^{-1}(0)=\num{124.835054800}$ \\
Scalar content & $N_H=2$ & $\bone=\frac{21}{5}$ with the same carrier family packet & $\alpha^{-1}(0)=\num{135.945991742}$ \\
External truncation & no $\nu^c$ & the family packet is externally reduced from $16$ to $15$ states, so $\Oadm=45$ and the hard carrier theorem no longer applies unchanged & $\alpha^{-1}(0)=\num{137.033841585}$ \\
\bottomrule
\end{tabularx}
\caption{Discrete exclusion structure in the carrier reading. The ``no $\nu^c$'' row is an external truncation of the hard family packet, not a normal carrier alternative.}
\end{table}

In addition, a light seam-even colored triplet zero mode would require a light seam-even bosonic mode in $E_3$. The minimal carrier picture does not provide such a mode, so the colored partner is not on equal footing with the Higgs candidate in $E_2$.

\section{Relative Spectral Completion Targets}

\begin{tcolorbox}[colback=gray!5, colframe=gray!60, breakable,
    title={\textbf{Status: dynamical target, not part of the hard carrier core}}, fonttitle=\bfseries]
The relative spectral action is the natural place to unify the $\alpha$-, $\chi$-, and metric channels, but the heat-kernel completion is not yet theorem-hard. In v2.9 it should therefore appear as a completion target rather than as part of the rigid carrier packet.
\end{tcolorbox}

The natural completion ansatz is
\[
S_{\mathrm{rel}}
= \tr f(D/\chi)
- \tr f(D_{\mathrm{ref}}/\chi)
+ \frac{i\pi}{2}\Delta\eta_\Gamma
+ S_{\mathrm{defect}}.
\]
The three stationarity channels are then
\[
\frac{\delta S_{\mathrm{rel}}}{\delta\alpha}=0,
\qquad
\frac{\delta S_{\mathrm{rel}}}{\delta\chi}=0,
\qquad
\frac{\delta S_{\mathrm{rel}}}{\delta e}=0,
\]
with an independent $\omega$ variation left to a later Palatini/Einstein--Cartan completion.

\subsection{Why the denominator 192 pi squared is plausible}

For a Dirac-type Laplace operator written as
\[
D^2=-\nabla^2+\mathcal E,
\]
the local four-dimensional heat-kernel coefficient is
\[
a_2(D^2)=\frac{1}{16\pi^2}\int d^4x\sqrt{-g}\,
\tr\!\left(\mathcal E+\frac16 R\right).
\]
For a Dirac-type curvature coupling one has $\mathcal E\sim R/4+\cdots$, so the curvature part scales like
\[
\frac{1}{16\pi^2}\left(\frac14+\frac16\right)
\int d^4x\sqrt{-g}\,\tr(\mathbf 1)\,R
=
\frac{5}{192\pi^2}
\int d^4x\sqrt{-g}\,\tr(\mathbf 1)\,R.
\]
After relative admissible-minus-reference subtraction and TFPT normalization, an Einstein-channel coefficient with denominator $192\pi^2$ is therefore natural. What remains open is the exact numerator and the precise sector trace selected by the full relative operator.

\begin{conjecture}[Induced Einstein channel]
The relative spectral completion should organize the metric channel as
\[
S_{\mathrm{EH}}^{\mathrm{rel}}
=
\frac{\chi^2}{192\pi^2}\,\tr_{\Hch}\mathbf 1
\int d^4x\sqrt{-g}\,R.
\]
On the canonical occupancy branch, where $\tr_{\Hch}\mathbf 1=\Oadm=48$, this gives
\[
S_{\mathrm{EH}}^{2.9}
= \int d^4x\sqrt{-g}\,\frac{\chi^2}{4\pi^2}R
= \int d^4x\sqrt{-g}\,\frac{\Mpl^2}{2}R,
\]
so that
\[
\Mpl^2=\frac{\chi_0^2}{2\pi^2}.
\]
\end{conjecture}

Using the observed reduced Planck mass only as a calibration check gives
\[
\chi_0 = \sqrt2\,\pi\,\Mpl \approx \num{1.08184e19}\ \mathrm{GeV}.
\]
The v2.9 claim is not that this is already hard, but that the relative-$a_2$ channel is the correct place to make it hard.

\section{Downstream Reinterpretations: \texorpdfstring{E$_8$}{E8}, Higgs, Flavor}

\subsection{\texorpdfstring{E$_8$}{E8} as a downstream carrier rereading}

The strongest E$_8$ rereading to retain is
\[
\lambda_{E8}=\frac{\gamma}{\ln(248/60)},
\qquad
\gamma=\tr_EY^2=\frac56.
\]
The E$_8$ sector is downstream of the carrier norm, not part of the hard minimal carrier core. One suggestive but non-rigid echo is $60=48+12=\Oadm+\dim\mathfrak g_{\SM}$.\footnote{This identity is numerically suggestive, but not theorem-level and should remain secondary or appendix-level.}

\subsection{Higgs as a conditional seam-even mode in \texorpdfstring{$E_2$}{E2}}

If a light seam-even bosonic zero mode exists in $E_2$, then the minimal carrier provides a natural Higgs candidate
\[
\tau H = +H,
\qquad
H\in \Gamma(E_2)\cong (1,2)_{1/2}.
\]
This is the correct representation-theoretic slot for a Standard Model Higgs doublet. The existence of the light mode itself remains part of the dynamical completion and should not be overstated as a theorem of the hard carrier core.

